\section{Risk Aversion and Implementation}

Risk aversion matters because a scoring rule can be implemented in different ways.
The distinction is between paying the score directly as money and using the score to determine the probability of a fixed prize.

\paragraph{Direct monetary scores.}
If the subject is paid the score itself, the subject solves
\[
\argmax_{r\in\OmegaN}\E_p[u(S(r,\omega))].
\]
For the frequency-guessing benchmark, this does not change incentives.
The payment is a fixed prize \(B\) if \(r=\omega\) and zero otherwise, so
\[
\E_p[u]
=
u(0)+\Pr_p(\omega=r)\{u(B)-u(0)\}.
\]
Any strictly increasing \(u\) therefore preserves the same optimal reports as maximizing the exact-match probability.

Direct monetary distance scores do not have this robustness in general.
Quadratic and Manhattan distance generate several possible payment levels, so nonlinear utility can change the ranking of reports.
This is the standard risk-aversion problem for monetary scoring rules \cite{ArmantierTreich2012,SchlagEtAl2013}.

\paragraph{Probability implementation.}
The distance rules can instead be implemented as probabilities.
Let \(q(r,\omega)\in[0,1]\) be a normalized score, and pay a fixed prize \(B\) with probability \(q(r,\omega)\), with zero otherwise.
Assume the prize and outside payment are fixed, utility depends only on final money, the randomization is objective and independent, and \(q\) is a valid probability.
Then
\[
\E_p\left[
q(r,\omega)u(B)+(1-q(r,\omega))u(0)
\right]
=
u(0)+(u(B)-u(0))\E_p[q(r,\omega)].
\]
Thus maximizing expected utility is equivalent to maximizing the expected normalized score, independently of the curvature of \(u\).

Because \(\OmegaN\) is finite, each rule considered here is bounded.
A positive affine normalization of a bounded score into \([0,1]\) preserves the expected-score ranking of reports.
Non-affine transformations generally need not preserve the report region.
This is the standard binary-lottery or binarized-scoring-rule logic used to separate incentive rankings from utility curvature \cite{ArmantierTreich2012,SchlagEtAl2013}.
The point is practical rather than novel: direct monetary distance rules are risk-neutral mechanisms, while probability implementations can make them robust to risk aversion under expected utility.
Thus the design comparison above should be read jointly with the implementation choice.
A researcher who values robustness may prefer fixed-prize discrete-metric scoring, while a researcher who wants quadratic-distance belief bounds can use a probability implementation to avoid relying on risk neutrality.
